
%%FILE blueprint/src/content.tex
% In this file you should put the actual content of the blueprint.
% It will be used both by the web and the print version.
% It should *not* include the \begin{document}
%
% If you want to split the blueprint content into several files then
% the current file can be a simple sequence of \input. Otherwise It
% can start with a \section or \chapter for instance.

\input{padic_fields}

\input{hilbert_symbol}

\input{quadratic_forms}

\printbibliography

%%FILE blueprint/src/hilbert_symbol.tex
\chapter{Hilbert symbol}

\section{Hilbert symbol over a field}

\begin{definition}
	\label{def:Hilbert_symbol}
	\lean{hilbertSym}
	\leanok
	Let $k$ be a field and $a, b \in k$. If both $a,b$ are nonzero, the \emph{Hilbert symbol} $(a, b)_k$ is defined as
	follows: if there exist $x, y, z \in k$, not all zero, such that $z^2 = ax^2 + by^2$,
	then $(a, b)_k = 1$; otherwise, $(a, b)_k = -1$. If one of $a,b$ is zero, we set $(a, b)_k = 0$.
\end{definition}


\subsection{Basic properties of the Hilbert symbol}

The Hilbert symbol has the following properties, useful for computing it in practice.

\begin{lemma}
	\label{lem:well_defined_up_to_squares}
	\lean{hilbertSym.mul_square_eq}
	\uses{def:Hilbert_symbol}
	\leanok
Let $a, b, a', b' \in k^\times$. Then $(a a'^2, b b'^2)_k = (a,b)_k$.
\end{lemma}

\begin{proof}
	\leanok
	Let $x, y, z \in k$ be such that $z^2 = a a'^2 x^2 + b b'^2 y^2$. Then, $xa', yb', z$ are not all
	zero and satisfy $z^2 = ax^2 + by^2$. Conversely, if $x, y, z \in k$ are such that
	$z^2 = ax^2 + by^2$, then $x/a', y/b', z$ are not all zero and satisfy
	$z^2 = a a'^2 x^2 + b b'^2 y^2$.
\end{proof}

The following lemma uses the definition \code{QuadraticAlgebra} from \mathlib.

\begin{lemma}
	\label{lem:eq_one_iff}
	\lean{hilbertSym.eq_one_iff}
	\uses{def:Hilbert_symbol}
	\leanok
	Let $a, b \in k^\times$. Then $(a, b)_k = 1$ if and only if there exists $t \in k(\sqrt{b})$ such
	that $a = N_{k(\sqrt{b})/k}(t)$.
\end{lemma}

\begin{proof}
	\leanok
	\begin{itemize}
		\item Case $b$ is a square in $k$. Then $k(\sqrt{b}) = k$, and the norm map is the identity, so
		$a$ is always the norm of itself. Let $c \in k^\times$ be such that $c^2 = b$.
		Then $c^2 = a \cdot 0^2 + b \cdot 1^2$, so $(a,b)_k=1$.
		\item Case $b$ is not a square in $k$. Then $k(\sqrt{b})$ is a quadratic extension of $k$, and
		the norm map is given by $N_{k(\sqrt{b})/k}(p + q \sqrt{b}) = p^2 - b q^2$.
		Suppose $a$ is a norm, so it satisfies $a = p^2 - b q^2$ for some $p, q \in k$.
		Then $1, q, p$ are not all zero and satisfy $ p^2 = a \cdot 1^2 + b \cdot q^2$,
		so $(a, b)_k = 1$. Conversely, if $x, y, z$ are not all zero and satisfy $z^2 = ax^2 + by^2$,
		then $x=0$, since otherwise $b$ would be a square in $k$, which contradicts our assumption.
		Then $x \ne 0$, and we have $a = (z/x)^2 - b (y/x)^2$, so $a$ is a norm.
	\end{itemize}
\end{proof}

\begin{proposition}
	\label{prop:hilbert_symbol_properties}
	\lean{hilbertSym.comm, hilbertSym.right_square_eq_one, hilbertSym.right_neg_self_eq_one,
	hilbertSym.right_one_minus_self_eq_one, hilbertSym.right_mul_eq_of_eq_one , hilbertSym.right_neg_mul,
	hilbertSym.right_minus_self_mul}
	\uses{def:Hilbert_symbol}
	\leanok
	The Hilbert symbol statisfies the following formulas:
	\begin{enumerate}[(i)]
		\item For $a, b \in k$, we have $(a, b)_k = (b, a)_k$.
		\item Let $a, c \in k^\times$. Then $(a, c^2)_k = 1$.
		\item For $a \in k^\times$, we have $(a, -a)_k = 1$.
		\item Let $a \in k$ with $a \ne 0, 1$. Then $(a, 1 - a)_k = 1$.
		\item Let $a, a', b \in k^\times$ and assume that $(a, b)_k = 1$. Then $(a a', b)_k = (a', b)_k$.
		\item Let $a, b \in k$. Then $(a, -ab)_k = (a, b)_k$.
		\item Let $a, b \in k^\times$ with $a \ne 1$. Then $(a, (1-a)b)_k = (a, b)_k$.
	\end{enumerate}
\end{proposition}

\begin{proof}
	\leanok
	\uses{def:Hilbert_symbol, lem:eq_one_iff, lem:well_defined_up_to_squares}

	\begin{itemize}
		\item Trivial if $a$ or $b$ are zero. Otherwise, if $x,y,z$ is a solution to $z^2 = ax^2 + by^2$, then $y,x,z$ is a solution to
	$z^2 = bx^2 + ay^2$.
		\item In this case, $c^2 = a \cdot 0^2 + b \cdot 1^2$, so $(a,b)_k=1$.
		N.B. Alternatively: the same argument as in the first case of the proof of Lemma~\ref{lem:eq_one_iff} applies.
		\item In this case, $0 = a \cdot 1^2 + (-a) \cdot 1^2$.
		\item In this case, $1 = a \cdot 1^2 + (1 - a) \cdot 1^2$.
		\item We know, by Lemma~\ref{lem:eq_one_iff}, that there exists $t \in k(\sqrt{b})$ such that
		$a = N_{k(\sqrt{b})/k}(t)$. If $(a',b)_k=1$, then $a'$ is the norm of some $t' \in k(\sqrt{b})$,
		so $aa' = N_{k(\sqrt{b})/k}(t t')$, so $(aa', b)_k = 1$. Conversely, if $(aa', b)_k = 1$, then
		$aa' = N_{k(\sqrt{b})/k}(t'')$ for some $t'' \in k(\sqrt{b})$, so
		$a' = N_{k(\sqrt{b})/k}(t''/t)$, hence $(a', b)_k = 1$.
		\item Trivial if $a$ or $b$ are zero. Otherwise it follows from the previous steps, since $(a, -a)_k = 1$.
		\item Follows from the previous steps, since $(a, 1 - a)_k = 1$.
	\end{itemize}
\end{proof}
\subsection{Local Properties}
We now specialize to the case where $k$ is either $\R$ or $\Q_p$ for some prime $p$. In this case, the Hilbert symbol can be computed explicitly.

\begin{theorem}
	\label{thm:real}
	\lean{hilbertSym.real_eq}
	\uses{def:Hilbert_symbol}
	\leanok
	Let $a, b \in \R^\times$. Then $(a, b)_{\R} = 1$ if and only if $a > 0$ or $b > 0$; otherwise $(a, b)_{\R} = -1$.
\end{theorem}

\begin{proof}
	If $a>0$, then $(1,0,\sqrt{a})$ works. Similarly, if $b>0$, then $(0,1,\sqrt{b})$ works.
	If both $a$ and $b$ are negative, then $ax^2 + by^2 \le 0$ for all $x,y$, so there is no
	nontrivial solution to $z^2 = ax^2 + by^2$.
\end{proof}

The following uses the \mathlib API for $p$-adic numbers and integers, as well as the results from the files \code{HassePrinciple/Padics/Lemmas.lean} and \code{HassePrinciple/Padics/Legendre.lean}.


For the statement of the theorem in the $p$-adic case, we need the following definitions.


\begin{definition}
	\label{def:epsilon}
	\lean{PadicInt.epsilon}
	\leanok
	Let $u \in (\Z_2)^\times$. We define $\epsilon(u)$ to be the class modulo $2$ of $\dfrac{u-1}{2}$.
\end{definition}

\begin{definition}
	\label{def:omega}
	\lean{PadicInt.omega}
	\leanok
	Let $u \in (\Z_2)^\times$. We define $\omega(u)$ to be the class modulo $2$ of $\dfrac{u^2-1}{8}$.
\end{definition}



For the proof, we will need the following lemma.
\begin{lemma}
	\label{lem:nontrivialzero}
	\lean{Padic.exists_nontrivial_zero}
	\leanok
Let $p$ be a prime and let $v \in (\Q_p)^\times$. Let $x, y, z \in \Q_p$ be such that
$(x, y, z) \ne (0, 0, 0)$ and $z^2 - p x^2 - v y^2 = 0$. Then there exist
$z', y' \in (\Z_p)^\times$ and $x' \in \Z_p$ such that $(z', y', x')$ is a nontrivial solution to
the same equation.
\end{lemma}
\begin{proof}
	Let $(x,y,z)$ be any nontrivial solution. For any $\lambda \ne 0$, also $(\lambda x, \lambda y,
	\lambda z)$ is a solution. Taking $\lambda = p^a$ for an appropriate $a$, we can assume that
	$x,y,z$ are all in $\Z_p$ and at least one of them is a unit. Now, if $z$ is not a unit, then
	$px^2+vy^2$ is also not a unit, thus $y$ is not a unit. Conversely if $y$ is not a unit, then
	$z^2-px^2$ is also not a unit, thus $z$ is not a unit. This shows that if either $y$ or $z$ is not
	a unit, then they are both non units, but then $p^2$ would divide $px^2$, hence $x$ would not be a
	unit either, contradicting the fact that one of $x,y,z$ is a unit.
\end{proof}

We are now ready to compute the Hilbert symbol $(a,b)$ in the $p$-adic case.
We split the statement and the proof of the theorem for the $p$-adic into the cases where $p$ is odd
and where $p = 2$, and for each we first consider separately the cases where the valuations of $a$
and $b$ are both $0$, both $1$, or one of them is $0$ and the other is $1$.

Everywhere in the following, we assume that $a, b \in (\Q_p)^\times$, and denote by $v_a$ and $v_b$
their valuations, and by $u_a$ and $u_b$ their unit parts.

\begin{theorem}
	\label{thm:padic_odd_p}
	\lean{hilbertSym.padic_odd_eq, hilbertSym.padic_odd_case00, hilbertSym.padic_odd_case10, hilbertSym.padic_odd_case11}
	\uses{def:Hilbert_symbol, def:padicUnitPart, def:epsilon, def:legendreSym}
	\leanok
Let $p$ be an odd prime and let $a, b \in (\Q_p)^\times$. Then
\[(a,b) = (-1)^{v_a v_b\epsilon(p)}\left(\dfrac{u_a}{p}\right) \left(\dfrac{u_b}{p}\right),\]
where $\left(\dfrac{\cdot}{p}\right)$ is the Legendre symbol.
\end{theorem}

\begin{proof}
	\uses{lem:nontrivialzero, prop:hilbert_symbol_properties, lem:multivariable_hensel}
	Since the Hilbert symbol is well defined up to squares, we can assume that $v_a, v_b \in \{0, 1\}$.
	\begin{itemize}
		\item Suppose $v_a=v_b=0$. Then $a=u_a$ and $b=u_b$ are units. By the theorem of Chevalley--Warning,
		the equation $z^2 = u_a x^2 + u_b y^2$ has a nontrivial solution in $\F_p$, so it has a
		nontrivial solution in $\Z_p$ (by multivariate Hensel's lemma), so $(a,b)=1$.
		Note: Chevalley--Warning is in \mathlib as \code{char_dvd_card_solutions}, in \code{Mathlib.FieldTheory.ChevalleyWarning}.
		\item Suppose $v_a=1$ and $v_b=0$. Then $a=p u_a$ and $b=u_b$. Since by the previous case, we
		have $(u_a, u_b)=1$, we have $(a,b) = (p u_a, u_b) = (p, u_b)$. We just need to show that this is
		equal to $\left(\dfrac{u_b}{p}\right)$. Observe that a unit is a square in $\Q_p$ if and only if
		it is a square modulo $p$. Now, if $u_b$ is a square, then both terms are $1$, see
		\ref{prop:hilbert_symbol_properties}. Otherwise, both terms are $-1$. In fact,
		if $z^2-p^2x^2-u_by^2$ has a nontrivial solution, we can assume that $z$ and $u_b$ are units, so
		we can reduce modulo $p$ to get $u_b = (zy^{-1})^2$, which is a contradiction.
		\item Suppose $v_a=v_b=1$. Then $a=p u_a$ and $b=p u_b$. This case follows from the properties of
		the Hilbert symbol and the previous case: $(pu_a,pu_b)=(pu_a, -p^2u_au_b)= (pu_a, -u_au_b)=
		\left(\dfrac{-u_au_b}{p}\right)$. By the properties of the Legendre symbol, this is equal to
		$\left(\dfrac{-1}{p}\right)\left(\dfrac{u_a}{p}\right)\left(\dfrac{u_b}{p}\right)$, as claimed.
	\end{itemize}
\end{proof}

We proceed similarly for the case $p = 2$.

\begin{theorem}
	\label{thm:padic_2}
	\lean{hilbertSym.two_adic_eq, hilbertSym.two_adic_case00, hilbertSym.two_adic_case10, hilbertSym.two_adic_case11}
	\uses{def:Hilbert_symbol, def:padicUnitPart, def:epsilon, def:omega, def:legendreSym}
	\leanok
Let $x,y \in (\Q_2)^\times$. Then
\[(a,b) = (-1)^{\epsilon(u_a)\epsilon(u_b) + v_a \omega(u_b) + v_b  \omega(u_a)}.\]
\end{theorem}

\begin{proof}
	\uses{prop:hilbert_symbol_properties, lem:nontrivialzero, lem:multivariable_hensel}
	Similarly to the previous case, we can assume that $v_a, v_b \in \{0, 1\}$.

	\noindent\textbf{Case 1} Suppose $v_a=v_b=0$. Then $a=u_a$ and $b=u_b$ are units. The claim is
		$(u_a,u_b)=(-1)^{\epsilon(u_a)\epsilon(u_b)}$, which is $1$ exactly when at least one of $u_a$
		and $u_b$ is congruent to $1$ modulo $4$ and $-1$ otherwise.
		Suppose first, without loss of generality, that $u_a \equiv 1 \pmod 4$. Then its residue class
		modulo $8$ is either $1$ or $5$. In the first case, $u_a$ is a square, so $(u_a,u_b)=1$. In the
		second case, $u_a+4u_b \equiv 1 \pmod 8$, so it is the square of some $w in \Z_2$, so
		\[
		w^2-u_a\cdot 1^2- u_b\cdot 2^2 = 0
		\]
		and thus $(u_a,u_b)=1$.

		If $u_a \equiv u_b \equiv -1 \pmod 4$, then suppose there exist $(x,y,z)$ not all zero, coprime,
		such that
		\[
		z^2-u_ax^2-u_by^2=0.
		\]
		Reducing modulo $4$, we obtain $z^2+x^2+y^2 \equiv 0 \pmod 4$, which is impossible, since the
		only squares modulo $4$ are $0$ and $1$, and the sum can be $0$ only if $x,y,z$ are all even,
		contradicting the fact that they are coprime. Thus, $(u_a,u_b)=-1$.
		This proof uses \code{ZMod.exists_sq_eq_two_iff} from \code{Mathlib.NumberTheory.LegendreSymbol.QuadraticReciprocity} and \code{hensels_lemma} in \code{Mathlib.NumberTheory.Padics.Hensel}

		\noindent\textbf{Case 2} Suppose $v_a=1$ and $v_b=0$. Then $a=2 u_a$ and $b=u_b$. We first show that
		\[
		(2, u_b) = (-1)^{\omega(u_b)},
		\]
		i.e. that $(2, u_b)=1$ if an only if $u_b \equiv \pm 1 \pmod 8$.
		Suppose that $(2, u_b)=1$. Equivalently, there exist $x,y,z$ not all zero, coprime, such that
		\[
		z^2-2x^2-u_by^2=0.
		\]
		By Lemma \ref{lem:nontrivialzero}, we can assume that $z$ and $y$ are units.
		Reducing modulo $8$, we get $1 - 2x^2 - u_b \equiv 0 \pmod 8$. Since squares modulo $8$ are
		$0, 1$ and $4$, we have $2x^2 \equiv 0$ or $2 \pmod 8$, so $u_b \equiv 1$ or $-1 \pmod 8$.

		Conversely, if $u_b \equiv 1 \pmod 8$, then it is a square, so $(2, u_b)=1$.
		If $u_b \equiv -1 \pmod 8$, then $z^2-2x^2-u_by^2=0$ has a solution modulo $8$, namely $(1,1,1)$,
		which can be lifted to a solution in $\Z_2$ by Hensel's lemma (\ref{lem:multivariable_hensel}),
		so $(2, u_b)=1$. Now, if we prove that $(2u_a,u_b)=(u_a,u_b)(2,u_b)$, then we are done. If either
		factor is $1$, by Proposition \ref{prop:hilbert_symbol_properties}, we are done. Otherwise, both
		factors are $-1$, so by the above we know $u_a \equiv -1 \pmod 4$ and $u_b \equiv 3 \pmod 8$. Up
		to multiplying by squares, we can consider the following two cases:
		\begin{itemize}
			\item $u_a=-1,u_b=3$. Then $(1,1,1)$ is a nontrivial solution to $z^2+2x^2-3y^2=0$.
			\item $u_a=3,u_b=-5$. Then $(1,1,1)$ is a nontrivial solution to $z^2-6x^2+5y^2=0$.
		\end{itemize}

		\noindent\textbf{Case 3} Suppose $v_a=v_b=1$. Using Proposition \ref{prop:hilbert_symbol_properties} and the proof
		of the previous case, we have
		\[
		(2u_a,2u_b)=(2u_a,-4u_au_b)=(2u_a,-u_au_b)=(-1)^{\epsilon(u_a)\epsilon(-u_au_b)+\omega(-u_au_b)},
		\]
		and a direct calculation shows that
		\[
		\epsilon(u_a)\epsilon(-u_au_b)+\omega(-u_au_b) = \epsilon(u_a)\epsilon(u_b) + \omega(u_b) + \omega(u_a).
		\]
\end{proof}

\begin{theorem}
		\label{thm:bilinear}
		\lean{Real.instHasBilinHilbertSym, Padic.instHasBilinHilbertSym}
		\uses{def:Hilbert_symbol}
		\leanok
The Hilbert Symbol is bilinear in both variables, i.e. for all $a, a', b, b' \in k^\times$, we have
\[
(aa',b)=(a,b)(a'b) \quad \text{and} \quad (a,bb')=(a,b)(a,b').
\]
\end{theorem}


\begin{proof}
\uses{thm:padic_odd_p, thm:padic_2}
By symmetry, it is enough to show the first equality.
If any of $a, a', b$ is zero, then both sides are zero. We will assume that $a,a',b \neq 0$.

The real case is trivial: if $b > 0$, then all the Hilbert symbols are $1$; otherwise,
if $a$ and $a'$ have the same sign, then $(aa',b) = 1$ because $aa' > 0$, and $(a,b) = (a',b)$, so
the right-hand side is also $1$;
if $a$ and $a'$ have different signs, then $(aa',b) = -1$ because $aa' < 0$, and $(a,b) = -(a',b)$,
so the right-hand side is also $-1$.

We now assume that $k=\Q_p$. Suppose first that $p$ is odd. By the formula in Theorem \ref{thm:padic_odd_p},
\begin{align*}
(aa,b) = (-1)^{v_{aa'}v_b\epsilon(p)}\left(\dfrac{u_{aa'}}{p}\right)^{v_b}\left(\dfrac{u_b}{p}\right)^{v_{aa'}}=\\
=(-1)^{v_a+v_{a'}v_b\epsilon(p)}\left(\left(\dfrac{u_{a}}{p}\right)\left(\dfrac{u_{a'}}{p}\right)\right)^{v_b}\left(\dfrac{u_b}{p}\right)^{v_a+v_{a'}}=\\
=(-1)^{v_av_b\epsilon(p)}(-1)^{v_{a'}v_b\epsilon(p)}\left(\dfrac{u_{a}}{p}\right)^{v_b}\left(\dfrac{u_{a'}}{p}\right)^{v_b}\left(\dfrac{u_b}{p}\right)^{v_a}\left(\dfrac{u_b}{p}\right)^{v_{a'}}=\\
(a,b)(a',b).
\end{align*}

If $p=2$, we use the formula in Theorem \ref{thm:padic_2}.
\begin{align*}
(aa',b) = (-1)^{\epsilon(u_{aa'})\epsilon(u_b) + v_{aa'} \omega(u_b) + v_b \omega(u_{aa'})} = (\bullet)
\end{align*}

Since $\epsilon$ and $\omega$ are group homomorphisms from a multiplicative group to an additive group, we have
\begin{align*}
	(\bullet) = (-1)^{(\epsilon(u_{a})+\epsilon(u_{a'}))\epsilon(u_b) + (v_{a}+v_{a'}) \omega(u_b) + v_b (\omega(u_a)+\omega(u_{a'}))}=\\
	=(-1)^{(\epsilon(u_{a})+\epsilon(u_{a'}))\epsilon(u_b)} (-1)^{(v_{a}+v_{a'}) \omega(u_b)} (-1)^{v_b (\omega(u_a)+\omega(u_{a'}))}=\\
	=(-1)^{\epsilon(u_{a})\epsilon(u_b)}(-1)^{\epsilon(u_{a'})\epsilon(u_b)} (-1)^{v_{a} \omega(u_b)}(-1)^{v_{a'} \omega(u_b)}(-1)^{v_b \omega(u_a)}(-1)^{v_b \omega(u_{a'})}=(a,b)(a',b)\\
\end{align*}

\end{proof}



\subsection{Global properties of the Hilbert symbol}
The main goal of this section is to prove
the product formula for the Hilbert symbol, which states that for every $a, b \in \Q^\times$, we have
\[\prod_{v \in V} (a, b)_v = 1\]
and the Existence Theorem, that states the existence of a rational number with prescribed Hilbert symbols.
Here, $V$ is the union of the set of prime numbers and the symbol $\infty$, and for each prime $p$,
the symbol $(a, b)_p$ is the Hilbert symbol computed in $\Q_p$, while $(a, b)_\infty$ is the Hilbert
symbol computed in $\R$.


\begin{theorem}
	\label{thm:almost_all_one}
	\lean{hilbertSym.almost_all_one}
	\uses{def:Hilbert_symbol}
	\leanok
Let $a, b \in \Q^\times$. Then $(a, b)_p = 1$ for all but finitely many primes $p$.
\end{theorem}

\begin{proof}
\leanok
	From the explicit formulas for the Hilbert symbol in the $p$-adic and in the real case, we see
	that it is multiplicative in both variables. 	In particular, it suffices to prove this theorem
	(and the following one) for $a$ and $b$ equal to $-1$ or
	a prime number. Using the formulas, we find:
	\begin{itemize}
	\item Case $a = -1, \, b = -1$. One has $(-1, -1)_\infty = (-1, -1)_2 = -1$ and
	$(-1, -1)_p = 1$ if $p \ne 2, \infty$;
  \item $a = -1, \, b = \ell$ with $\ell$ prime. If $\ell=2$, one has $(-1, 2)_v = 1$
	for all places $v$; if $v \ne 2$, one has $(-1, \ell)_v = 1$ if $v \ne 2, \ell$,
	and $(-1, \ell)_2 = (-1, \ell)_{\ell} = (-1)^{e(\ell)}$;
  \item $a = \ell, b = \ell'$ with $\ell, \ell'$ primes. If $\ell = \ell'$, by Proposition
	\ref{prop:hilbert_symbol_properties}, $(\ell, \ell)_v = (-1, \ell)_v$ for all $v$ and we are
	reduced to the case considered above. If $\ell \ne \ell'$ and if $\ell' = 2$, one has
	$(\ell, 2)_v = 1$ for $v \ne 2, \, \ell$ and
\[
(\ell, 2)_2 = (-1)^{\omega(\ell)}, \qquad (\ell, 2)_{\ell} = \left(\frac{2}{\ell}\right) = (-1)^{\omega(\ell)}, \quad \text{cf. chap. I, n° 3.2, th. 5.}
\]
If \( \ell \) and \( \ell' \) are distinct and different from 2, one has \( (\ell, \ell')_v = 1 \) for \( v \neq 2, \, \ell, \, \ell' \) and
\[
(\ell, \ell')_2 = (-1)^{\epsilon(\ell)e(\ell')}, \quad (\ell, \ell')_{\ell} = \left(\frac{\ell'}{\ell}\right), \quad (\ell, \ell')_{\ell'} = \left(\frac{\ell}{\ell'}\right).
\]
\end{itemize}
\end{proof}

\begin{theorem}
	\label{thm:product_formula}
	\lean{hilbertSym.prod_eq_one}
	\uses{def:Hilbert_symbol}
	\leanok
Let $a, b \in \Q^\times$. Then $\prod_{v \in V} (a, b)_v = 1$.
\end{theorem}
\begin{proof}
	\uses{thm:almost_all_one, thm:real, thm:padic_odd_p, thm:padic_2, prop:hilbert_symbol_properties}
	Follows immediately from the proof of Theorem \ref{thm:almost_all_one} and quadratic reciprocity,
	which can be imported from \code{Mathlib.NumberTheory.LegendreSymbol.QuadraticReciprocity}.
	In all cases, the product is one.
\end{proof}

\subsection{Approximation and Existence Theorem}

\begin{definition}
	\label{def:finite_embedding}
	\lean{Rat.finiteEmbedding}
	\leanok
Let $S$ be a finite set of prime numbers.
We define the \emph{finite embedding} of $\Q$ into $\R \times \prod_{p \in S} \Q_p$ as the diagonal map.

\end{definition}

Note: we will care about the product of finitely many completions of $\Q$, which do not necessarily
include $\R$, but the proof of the Approximation Theorem assumes it without loss of generality,
since we can always add $\infty$ to $S$.

\begin{theorem}\label{thm:approximation}%todo change references
	\lean{Rat.approximation}
	\uses{def:finite_embedding}
	\leanok
	Let $S$ be a finite set of prime numbers.
	The image of the finite embedding of $\Q$ into $\R \times \prod_{p \in S} \Q_p$ is dense with respect to
	the product topology. More concretely, for every $\epsilon > 0$ and every
	$y \in \R \times \prod_{p \in S} \Q_p$, there exists $x \in \Q$ such that the distance
	between $y$ and the image of $x$ is less than $\epsilon$.
\end{theorem}
\begin{proof}
	Let $y=(y_\infty, (y_p)_{p \in S})$ be an element of $\R \times \prod_{p \in S} \Q_p$.
	We can assume $y_p \in \Z_p$ for all $p \in S$, since otherwise we can multiply all $y_p$'s by an
	integer. Let $\varepsilon > 0$ and $N \ge 0$. We need to prove that there exists $x \in \Q$ such that
	$|y_\infty - x| \le \varepsilon$ and $v_p(y_p - x) \ge N$ for all $p \in S$.

	By the Chinese Remainder Theorem \code{Nat.chineseRemainderOfFinset}, there exists $x_0 \in \N$
	such that $x_0 \equiv y_p \pmod{p^N}$ for all $p \in S$.

	Fix $q$ a prime number. By density of the numbers of the form $a/q^m$ in $\R$, where $a \in \Z$
	and $m \ge 0$, we can find $a$ in $\Z$, $q$ a prime not in $S$ and $m \ge 0$, such that
	\[
	|x_0 - x_\infty - \frac{a}{q^m}(p_1\dots p_n)^N| \le \varepsilon.
	\]

	Then $x = x_0 + \frac{a}{q^m}(p_1\dots p_n)^N$ satisfies the required conditions.
\end{proof}


\begin{theorem}\label{thm:existence}
	\lean{hilbertSym.exists_rat_with_finite_prescribed_hilbertSym}
	\uses{def:Hilbert_symbol}
	\leanok
	Let $(a_i)_{i \in I}$ be a finite family of nonzero rational numbers, and let
	$e_{i,v} \in \{\pm 1\}$ for each $i \in I$ and each place $v$ of $\Q$. There exists a rational
	number $x \in \Q^\times$ such that
	\[
	\forall i \in I, \forall v \text{ place of } \Q, (x,a_i)_v=e_{i,v}
	\]
	if and only if the following conditions hold:
	\begin{enumerate}
		\item For each $i \in I$, we have $e_{i,v}=1$ for all but finitely many places $v$.
		\item For each $i \in I$, we have $\prod_v e_{i,v}=1$.
		\item For each place $v$, there exists $x_v \in \Q_v^\times$ such that
		$(x_v, a_i)_v = e_{i,v}$ for all $i \in I$.
	\end{enumerate}
\end{theorem}

The proof uses \code{Nat.infinite_setOf_prime_and_modEq}.

\begin{proof}
	\uses{def:Hilbert_symbol, thm:approximation, thm:real, thm:padic_odd_p, thm:padic_2,
	prop:hilbert_symbol_properties, thm:product_formula, thm:padic_squares_open}
To show that the conditions are necesary, simply observe that 1. is Theorem \ref{thm:almost_all_one},
2. is Theorem \ref{thm:product_formula}, and 3. is trivial since we can take $x_v = x$ for all $v$.

Conversely, suppose that 1., 2.\ and 3.\ hold. Up to multiplying by a square integer (which does
not change the Hilbert symbols), we can assume that all $a_i$'s are integers.
Let $S$ be the set of places $2$, $\infty$ and all the prime divisors of the $a_i$'s.
Let $T$ be the set of places $v$ such that there exists $i \in I$ with $e_{i,v}=-1$.
The set $S$ is finite and, by 1., $T$ is also finite.

Consider first the case where $S \cap T= \emptyset$. Then, we define
\[
a = \prod_{\ell \in T} \ell, \quad m = 8 \prod_{\ell \in S \setminus\{2,\infty\}} \ell.
\]

Since $a,m$ are coprime, by \code{Nat.infinite_setOf_prime_and_modEq}, there exists a prime number
$p$ such that $p \equiv a \pmod m$ and $p \not\in S \cup T$.

Taking $x=ap$ works.
\begin{itemize}
	\item If $v= \infty$, then $x > 0$, so $(x, a_i)_\infty = 1$ for all $i$, and $\infty \notin T$,
	so $e_{i,\infty}=1$.
	\item If $v = \ell \in S$, we have $x \equiv a^2 \pmod m$, so $x$ is a square modulo any $\ell \in S$,
	hence by Hensel's lemma it is a square in $\Q_\ell$, and $x$ is a square modulo $8$, so again by
	Hensel's lemma it is a square in $\Q_2$. By Proposition \ref{prop:hilbert_symbol_properties},
	$(x,a_i)=1$.
	\item If $v = \ell \notin T$ and $\ell \ne p$, $v_\ell(x)=0$ and so
	$(x,a_i)_{\ell}=\left(\dfrac{a_i}{\ell}\right)^v_{\ell}(x)=1$. On the other hand, $e_{i,\ell=1}$,
	since $v \notin T$.
	\item If $v = \ell \in T$, then $v_\ell(x)=1$, so $(x,a_i)_{\ell}=\left(\dfrac{a_i}{\ell}\right)$.
  By 3., there exists $x_\ell \in \Q_\ell^\times$ such that $(x_\ell, a_i)_\ell = e_{i,\ell}$ for all $i$,
	and $v_\ell(x_\ell)$ is odd since at least one of the $e_{i,\ell}$'s is $-1$.
	Therefore, $(a_i,x)_\ell$ and $(a_i,x_\ell)_\ell$ are both equal to $\left(\dfrac{a_i}{\ell}\right)$.
	\item If $v = p$, since for all the other places the equality holds, by the product formula we
	also have $(x,a_i)_p = e_{i,p}$ for all $i$.
\end{itemize}

For the general case, since the image of $\Q^\times$ is dense in $\prod_{v \in S} \Q_v^\times$, by
Theorem \ref{thm:approximation}, and squares are open in $\Q_v^\times$, by Theorem \ref{thm:padic_squares_open},
we can find $x' \in \Q^\times$ such that $x'/x_v$ is a square in $\Q_v^\times$ for all $v \in S$.
Then, $(x', a_i)_v = (x_v, a_i)_v = e_{i,v}$ for all $i$ and all $v \in S$.
Let $\eta_{i,v}= e_{i,v} (x', a_i)_v$. Then $\eta_{i,v}=1$ for all $v \in S$ and conditions 1., 2.\
and 3. hold for the $\eta_{i,v}$'s, so we can apply the first part of the proof to find $y \in \Q^\times$
such that $(y, a_i)_v = \eta_{i,v}$ for all $i$ and all $v$.

Finally, setting $x = x' y$, we have $(x, a_i)_v = (x', a_i)_v (y, a_i)_v = e_{i,v}$ (by the explicit
formulas for the Hilbert symbol in the $p$-adic and in the real case) for all $i$ and all $v$, as desired.
\end{proof}


%%FILE blueprint/src/padic_fields.tex
\chapter[p-adic Fields]{$p$-adic Fields}

\section[p-adic Equations]{$p$-adic Equations}\label{sec:padic_equations}

\begin{definition}\label{def:primitive}
\lean{Function.IsPrimitive}
\leanok
Let $R$ be a ring. An element $(x_1, \cdots, x_n) \in R^n$ is \emph{primitive} if at least one of the $x_i$ is invertible.
\end{definition}


\begin{proposition}\label{prop:common_root}
\lean{Function.IsPrimitive}
\leanok
\uses{def:primitive}
Let $f^{(i)} \in \Z_p[X_1, \cdots, X_m]$ be homogeneous polynomials with $p$-adic integer coefficients. For each $n \ge 1$, denote by $f_n^{(i)}$ the reduction of $f^{(i)}$ modulo $p^n$.
 The following are equivalent:
\begin{enumerate}[a)]
\lean{Padic.common_root_tfae}
\leanok
\item The $f^{(i)}$ have a nontrivial common zero in $\Q_p^m$.
\item The $f^{(i)}$ have a common primitive zero in $\Z_p^m$.
\item For all $n \ge 1$, the $f_n^{(i)}$ have a common primitive zero in $(\Z / p^n \Z_p)^m$.
\end{enumerate}
\end{proposition}

\begin{proof}
The implication (b) $ \implies$ (a) is trivial. Conversely, if $x = (x_1, \cdots, x_m) \in \Q_p^m$ is a nontrivial common zero of the $f^{(i)}$, put
\[ h := \min(v_p(x_1), \cdots, v_p(x_m)) \quad \text{and} \quad y := p^{-h} x. \]
Then $y \in \Z_p^m$ is a common primitive zero of the $f^{(i)}$.
For the equivalence (b) $\iff$ (c), Serre applies the first Lemma from \cite[\S II.2.1]{Ser73}. However, using \mathlib's \texttt{DirectLimit} definition here might be tricky. It could be easier to just prove this specific result; note that the file \texttt{Mathlib.NumberTheory.Padics.RingHoms} contains useful results relating $\Z_p$ and $\Z/p^n\Z$ (without relying on \texttt{DirectLimit}).

\end{proof}

\section[Squares in the p-adic units]{Squares in $\Q_p^{\times}$}\label{sec:padic_squares}

\begin{theorem}\label{thm:padic_squares_open}
\lean{Padic.unitSquares}
\leanok
${\Q_p^\times}^2$ is an open subgroup of $\Q_p^\times$.
\end{theorem}

\begin{proof}
\leanok
The only nontrivial part is to show that the squares form an open set.  Since $\Q_p$ is a field, the unit topology agrees with the subspace topology, so it suffices to check that the set of nonzero squares in $\Q_p$ is open. Define $r$ as $1$ if $p$ is odd, or as $2^{-2}$ if $p = 2$, and denote by $B_{1, r}$ the open ball with center $1$ and radius $r$. For each $c \in B_{1, r}$, by applying Hensel's lemma with the polynomial $X^2 - c$, we obtain that $c$ is a square.
Then, for each nonzero square $s \in \Q_p$, $s \cdot B_{1, r}$ is an open neighbourhood of $s$ contained in the set of nonzero squares, which shows that the latter is open.
See \cite[\S II.3.3]{Ser73} for more details.
\end{proof}

\begin{lemma}
  \label{lem:multivariable_hensel}
  \lean{PadicInt.multivariable_hensel}
  \leanok
Let $p$ be a prime, let $f \in \Z_p[X_1, \cdots, X_m]$ be a polynomial with $p$-adic integer
coefficients, and let $x \in (\Z/p\Z)^m$ be a zero of $f$ modulo $p^n$, for some positive integer
$n$. Suppose that, for some $j \in \{1, \cdots, m\}$, the partial derivative
$\frac{\partial f}{\partial X_j}(x)$ is nonzero modulo $p^k$, where $0 < 2k < n$.
Then there exists a zero $z$ of $f$ in $\Z_p^m$ whose reduction modulo $p^n$ is $x$, and such that
$z-x$ is divisible by $p^{n-k}$.
\end{lemma}

\begin{proof}
This is in Serre's \cite[\S II.2 Theorem 1]{Ser73}.
\end{proof}

\section[Units and Legendre symbols]{Units and Legendre symbols}\label{sec:padic_legendre}

\begin{lemma}\label{lem:padic_unit_part_unit}
\lean{Padic.norm_mul_pow_neg_valuation_eq_one}
\leanok
Let $p$ be a prime, let $x \in \Q_p^\times$, and let $v_p$ denote the $p$-adic valuation. Then $\|x \cdot p^{-v_p(x)}\| = 1$.
\end{lemma}
\begin{proof}
\leanok
Straightforward.
\end{proof}

\begin{definition}\label{def:padicUnitPart}
  \lean{Padic.unitPart}
  \uses{lem:padic_unit_part_unit}
  \leanok
Let $p$ be a prime and $x \in \Q_p^\times$. The \emph{unit part} of $x$ is defined as the element
of $\Z_p^\times$ corresponding to $x \cdot p^{-v_p(x)}$.
\end{definition}

\begin{definition}\label{def:legendreSym}
  \lean{PadicInt.legendreSym}
  \leanok
Let $p$ be a prime and $u \in \Z_p^\times$. The \emph{Legendre symbol} $\left(\dfrac{u}{p}\right)$
is defined as
\[\left(\dfrac{u}{p}\right)=\left(\dfrac{\overline{u}}{p}\right),\]
where $\overline{u}$ is the image of $u$ under the natural map $\Z_p^\times \to (\Z/p\Z)^\times$.
\end{definition}


%%FILE blueprint/src/quadratic_forms.tex
\chapter[Quadratic Forms over the Rationals and the p-adics]{Quadratic Forms over $\Q$ and $\Q_p$}

\section{Quadratic Forms}\label{sec:quadratic_form}

The following definition is available in \mathlib under the name \code{QuadraticForm}:
\begin{definition}\label{def:quadratic_form}
\mathlibok
Let $M$ be a module over a commutative ring $A$. A function $f : M \to A$ is called a \emph{quadratic form} on $M$ if
\begin{enumerate}
\item $f(ax) = a^2f(x)$ for all $a \in A, x \in M$,
\item The function $(x, y) \mapsto f(x + y) - f(x) - f(y)$ is a bilinear form.
\end{enumerate}
\end{definition}

The \code{QuadraticForm} folder of \mathlib contains definitions and results about quadratic forms, including the bilinear map
\[ (x, y) \mapsto x \cdot y := \frac{1}{2}(f(x + y) - f(x) - f(y))\] associated to a quadratic form (see \code{QuadraticMap.associated}) and the definition and characterization of nondegenerate quadratic forms (see \code{QuadraticMap.Nondegenerate}).

\begin{definition}\label{def:isotropic}
\lean{QuadraticMap.Isotropic}
\uses{def:quadratic_form}
\leanok
A quadratic form $f : M \to A$ is called \emph{isotropic} if there exists a nonzero $x \in M$ such that $f(x) = 0$.
\end{definition}

% NOTE: if needed, assume $R$ is a field.
\begin{definition}\label{def:represents}
	\lean{QuadraticMap.represents}
	\leanok
	Let $R$ be a commutative ring, let $V$ be an $R$-module and let $Q$ be a quadratic form on $V$. We say that $Q$ \emph{represents} an element $r \in R$ if there exists $x \in V$ such that $Q(v) = r$ and $v \ne 0$.
\end{definition}

\begin{lemma}\label{lem:represents_zero}
	\lean{QuadraticMap.represents_zero_iff_isotropic}
	\leanok
	\uses{def:represents, def:isotropic}
	Let $R$ be a commutative ring, let $V$ be an $R$-module and let $Q$ be a quadratic form on $V$. Then $Q$ \emph{represents} $0$ if and only if it is isotropic.
\end{lemma}
\begin{proof}
	\leanok
	This follows immediately from the definitions.
\end{proof}


\begin{lemma}\label{lem:represents_iff}
\lean{QuadraticMap.Equivalent.represents_iff, QuadraticMap.Equivalent.isotropic_iff}
\leanok
\uses{def:represents, def:isotropic, lem:represents_zero}
Let $Q$ and $Q'$ be two equivalent quadratic forms over $R$ and let $r \in $. Then $Q$ represents $r$ if and only if $Q'$ represents $r$. In particular, $Q$ is isotropic if and only if $Q'$ is isotropic.
\end{lemma}

\begin{proof}
\leanok
By Lemma \ref{lem:represents_zero}, it suffices to prove the first claim, which follows from the definition of equivalence of quadratic forms.
\end{proof}

\subsection{Orthogonality}\label{subsec:ortho}

Let $R$ be a commutative ring, let $V$ be an $R$-module and let $Q$ be a quadratic form on $V$.

\begin{definition}\label{def:ortho_compl}
\lean{QuadraticForm.orthoCompl}
\leanok
Let $S$ be a subset of $V$. The \emph{$Q$-orthogonal complement of $U$} is the $R$-submodule $U^o$ of $V$ consisting of all elements $v \in V$ that are $Q$-orthogonal to every $s \in S$.
\end{definition}

\begin{definition}\label{def:to_dual}
\lean{QuadraticForm.toDual}
\leanok
Let $U$ be an $R$-submodule of $V$. We define $q_U : V \to U^*$ as the $R$-linear map sending $v \in V$ to the linear form $(u \in U) \mapsto v \cdot u$.
\end{definition}

\begin{lemma}\label{lem:ker_to_dual}
Let $U$ be an $R$-submodule of $V$. Then the kernel of $q_U$ is $U^o$. In particular, the kernel of $q_V$ is equal to the radical of $Q$.
\lean{QuadraticForm.orthoCompl_eq_ker_toDual, QuadraticForm.radical_eq_orthoCompl_top}
\leanok
\uses{def:ortho_compl, def:to_dual}
\end{lemma}
\begin{proof}
\leanok
Follows by expanding the definitions.
\end{proof}

From now on, we assume that $R$ is a field.

\begin{lemma}\label{lem:nondegenerate_iff_to_dual_bijective}
\lean{QuadraticForm.nondegenerate_iff_toDual_bijective}
\uses{def:to_dual}
\leanok
$Q$ is nondegenerate if and only if $q_V$ is bijective.
\end{lemma}
\begin{proof}
\uses{lem:ker_to_dual}
\leanok
This follows from Lemma \ref{lem:ker_to_dual} and the fact that a quadratic form is nondegenerate if and only if its radical is trivial.
\end{proof}

\begin{lemma}\label{lem:to_dual_surjective}
Let $U$ be an $R$-vector subspace of $V$. If $Q$ is nondegenerate, then $q_U$ is surjective.
\lean{QuadraticForm.toDual_surjective}
\leanok
\uses{def:to_dual}
\end{lemma}
\begin{proof}
\uses{lem:nondegenerate_iff_to_dual_bijective}
\leanok
$q_U$ is equal to the composition of $q_V : V \to V^*$, which is surjective by Lemma \ref{lem:nondegenerate_iff_to_dual_bijective} because $Q$ is nondegenerate, and the canonical surjection $V^* \to U^*$.
\end{proof}

We need to introduce some notation before the next proposition. Let $S$ be a commutative ring, let $V$ and $W$ be $S$-modules and let $Q$ and $Q'$ be quadratic forms on $V$ and $W$, respectively. Then $Q \dotplus Q'$ denotes the quadratic form on $V \times W$ obtained by applying $Q$ on $V$ and $Q'$ on $W$. This is available in \mathlib as \texttt{QuadraticMap.prod}. Similarly, $Q \dotminus Q'$ denotes the form $Q \dotplus (-Q')$ on $V \times W$.

\begin{proposition}\label{prop:ortho_compl_properties}
\lean{QuadraticForm.finrank_eq_add, QuadraticForm.orthoCompl_orthoCompl, QuadraticForm.radical_restrict_eq_inf, QuadraticForm.nondegenerate_orthoCompl_iff, QuadraticForm.orthoCompl_isCompl, Submodule.prodOrthoComplEquiv}
\leanok
\uses{def:ortho_compl}
Let $U$ be an $R$-vector subspace of $V$ and suppose that $Q$ is nondegenerate. Denote by $Q_U$ the restriction of $Q$ to $U$. Then :
\begin{enumerate}[(i)]
\item $\dim U + \dim U^o = \dim V$.
\item $U^{oo} = U$.
\item $\text{rad}(U) = \text{rad}(U^o) = U \cap U^o$.
\item $Q_U$ is nondegenerate if and only if $Q_{U^o}$ is nondegenerate, in which case $V = U \oplus U^o$ and $Q$ is equivalent to $Q_U \dotplus Q_{U^o}$.
\end{enumerate}
\end{proposition}
\begin{proof}
\leanok
\uses{lem:to_dual_surjective}
See Proposition 2 in \S IV.1.2 of \cite{Ser73}.
\end{proof}

\subsection{Orthogonal basis}\label{subsec:ortho_basis}

Let $K$ be a field in which $2$ is invertible and let $V$ be a $K$-vector space.

The definition of orthogonal basis for $V$ is available in \mathlib as \code{LinearMap.IsOrtho}$_i$. The lemma \code{LinearMap.BilinForm.exists_orthogonal_basis} proves that for any symmetric bilinear form $B$ on $V$ there exists an orthogonal basis with respect to $B$.

\begin{definition}\label{def:contiguous}
\lean{Module.Basis.IsContiguous}
\leanok
Two bases of $V$ are called \emph{contiguous} if they have an element in common.
\end{definition}

\begin{theorem}\label{thm:exists_chain}
\lean{Module.Basis.chainOfNondegenerate}
\uses{def:chain, lem:exists_const}
\leanok
Let $V$ be a $K$-vector space of dimension at least $3$, let $Q$ be a nondegenerate quadratic form on $V$, and let $b, b'$ be two $Q$-orthogonal bases of $V$. Then there exists a finite sequence $(b^{(0)}, \cdots, b^{(m)})$ of orthogonal bases of $V$ such that $b^{(0)} = b$, $b^{(m)} = b'$ and $b^{(i)}$ is contiguous with $b^{(i + 1)}$ for $0 \le i < m$.
\end{theorem}

\begin{proof}
\uses{def:contiguous}
See \cite[\S IV.1.4, Theorem 2]{Ser73}.
\end{proof}

\begin{definition}\label{def:chain}
\lean{Module.Basis.Chain}
\leanok
A sequence of bases $(b^{(0)} = b, \cdots, b^{(m)} = b')$ as in Theorem \ref{thm:exists_chain} is called a \emph{chain} of orthogonal bases contiguously relating $b$ to $b'$.
\end{definition}

\begin{lemma}\label{lem:exists_const}
\lean{Module.Basis.exists_const}
\leanok
Let $V$ be a $K$-vector space of dimension at least $3$, let $Q$ be a nondegenerate quadratic form on $V$, and let $b, b'$ be two $Q$-orthogonal bases of $V$. Assume that
\[ (b_1 \cdot b_1)(b_i' \cdot b_i') - (b_i' \cdot b_i')^2 = 0 \text{ for } i = 1, 2. \]
There exists $x \in K$ such that $b_x := b_1' + x b_2'$ is nonisotropic and generates with $b_1$ a nondegenerate plane.
\end{lemma}

\begin{proof}
	See \cite[\S IV.1.4, Lemma after Theorem 2]{Ser73}.
\end{proof}

\subsection{Translations}\label{subsec:translations}

\begin{definition}\label{def:hyperbolic}
\lean{QuadraticForm.IsHyperbolic}
\leanok
A quadratic form is \emph{hyperbolic} if it is equivalent to the form $XY$.
\end{definition}

\begin{lemma}\label{lem:represents_of_hyperbolic}
\lean{QuadraticForm.represents_of_isHyperbolic}
\leanok
\uses{def:hyperbolic, def:represents}
An hyperbolic quadratic form represents all elements of $R$.
\end{lemma}
\begin{proof}
\leanok
\uses{lem:represents_iff}
By Lemma \ref{lem:represents_iff}, it suffices to prove this for the quadratic form $XY$, which holds since for each $r \in R$, we have $XY(r, 1) = 0$.
\end{proof}


\begin{proposition}\label{prop:sum_hyperbolic}
\lean{QuadraticForm.equivalent_isHyperbolic_add, QuadraticForm.represents_of_isotropic_of_nondegenerate}
\leanok
\uses{def:represents, def:hyperbolic}
Let $K$ be a field. If a quadratic form $Q$ on an $K$-vector space $V$ represents $0$ and is nondegenerate, then $Q$ is equivalent to $H + Q'$, where $H$ is hyperbolic. Moreover, $f$ represents all elements of $K$.
\end{proposition}
\begin{proof}
\leanok
\uses{lem:represents_of_hyperbolic, prop:ortho_compl_properties}
Fix $x \in V$ such that $Q(x) = 0$. Since $Q$ is nondegenerate, there exists $z \in V$ such that $x \cdot z = 1$. The element $y = z - 1/2 \cdot(z\cdot z)x$ is isotropic and $x \cdot y = 1$.
The restriction of $Q$ to the submodule $U := Kx \oplus Ky$ is the desired hyperbolic form $H$. Take $Q'$ to be the restriction of $Q'$ to the orthogonal complement of $U$; then $Q$ is equivalent to $H + Q'$ by Proposition \ref{prop:ortho_compl_properties}(iv).
Since $H$ represents all elements of $K$ by Lemma \ref{lem:represents_of_hyperbolic}, so does $Q = H + Q'$.
\end{proof}

\begin{corollary}\label{cor:represents_iff_sub_sub}
\lean{QuadraticForm.represents_iff_sub_isotropic}
\leanok
Let $K$ be a field, let $r \in K^\times$ and let $Q$ be a quadratic form on a free $K$-vector space $V$. Then $Q$ represents $r$ if and only if the form $Q' := Q \dotminus rZ$ represents $0$.
\end{corollary}
\begin{proof}
\uses{prop:sum_hyperbolic}
The forward implication is clear. For the second one, fix a basis for $V$ and assume that $Q' := Q \dotminus rZ$ has a nontrivial zero $(x_1, \cdots, x_n, z) \in V \times K$. Then either $z = 0$, in which case $Q$ represents zero and thus also $r$, or $Q(x_1/z, \cdots, x_n/z) = r$.
\end{proof}

\begin{corollary}\label{cor:prod_isotropic_iff}
\lean{QuadraticForm.prod_isotropic_iff, QuadraticForm.prod_isotropic_iff'}
\leanok
Let $K$ be a field and $Q$ and $Q'$ be nondegenerate quadratic forms over $K$ of rank at least $1$. Then the following properties are equivalent:
\begin{enumerate}[a)]
\item $Q \dotminus Q'$ represents $0$.
\item There exists $r \in K^\times$ which is represented by $Q$ and $Q'$.
\item There exists $r \in K^\times$ such that $Q \dotminus rZ^2$ and $Q' \dotminus rZ^2$ represent $0$.
\end{enumerate}
\end{corollary}

\begin{proof}
\uses{cor:represents_iff_sub_sub}
The equivalence (b) $\iff$ (c) follows from Corollary \ref{cor:represents_iff_sub_sub}. The implication (b) $\implies$ (a) is trivial. To show (a) $\implies (b)$, write a nontrivial zero of $Q \dotminus Q'$ in the form $(x, y)$ with $Q(x) = Q'(y)$. If $Q(x) \ne 0$, we are done. If $Q(x) = 0$, then $Q$ represents $0$ and thus by Proposition \ref{prop:sum_hyperbolic} also all elements of $K$. In particular, it represents all nonzero values taken by $Q'$.
\end{proof}

\section[Quadratic Forms over the p-adics]{Quadratic Forms over $\Q_p$}\label{sec:quadratic_form_padic}

\begin{definition}\label{def:discr}
\mathlibok
\lean{QuadraticForm.discr}
Let $p$ be a prime number and let $f$ be a quadratic form of rank $n$ over $\Q_p$. Take a diagonal quadratic form $a_1X_1^2 + \cdots + a_nX_n^2$ which is equivalent to $f$; then the \emph{discriminant} of $f$ is the element $d(f) = a_1 \cdots a_n \in \Q_p^\times/{\Q_p^\times}^2$.

Note that \mathlib contains \code{QuadraticForm.discr}, as an element of $\Q_p$. We might need to adjust some statements and proofs if we really need to regard it as an element in $\Q_p^\times/{\Q_p^\times}^2$.
\end{definition}

\begin{definition}\label{def:HM_inv}
\uses{def:Hilbert_symbol}
\lean{QuadraticForm.hasseMinkowskiInv}
\leanok
Let $p$ be a prime number and let $f$ be a quadratic form of rank $n$ over $\Q_p$. Take a diagonal quadratic form $a_1X_1^2 + \cdots + a_nX_n^2$ which is equivalent to $f$; then the \emph{Hasse--Minkowski invariant} of $f$ is the element
\[ \epsilon(f) := \prod_{i < j}(a_i, a_j).\]
\end{definition}

\begin{lemma}\label{lem:HM_inv_well_defined}
\uses{def:HM_inv, thm:exists_chain}
\lean{QuadraticForm.hasseMinkowskiInv.eq_of_equivalent}
\leanok
The value $\epsilon(f)$ does not depend on the choice of diagonal quadratic form equivalent to $f$.
\end{lemma}
\begin{proof}
Let $b = (b_1, \cdots, b_n)$ be an orthogonal basis for $V$ such that $Q(x) = a_1 x_1^2 + \cdots + a_n x_n^2$ for $x = \sum x_i b_i$. We denote  $\epsilon(b) := \prod_{i < j}(a_i, a_j)$.

If $n=1$, then $\epsilon(f) = 1$. If $n=2$, then $\epsilon(f) = 1$ if and only if the form $Z^2 - a_1 X^2 - a_2 Y^2$ represents $0$. By Corollary \ref{cor:represents_iff_sub_sub}, this is equivalent to $a_1 X^2 + a_2 Y^2$ representing $1$, which means that there exists $v \in V$ with $Q(v) = 1$; this is independent on $b$.
For $n \ge 3$ we use induction on $n$. By Theorem \ref{thm:exists_chain}, it suffices to prove that for any pair of contiguous bases $b$ and $b'$, we have $\epsilon(b) = \epsilon(b')$.
Thanks to the symmetry of the Hilbert symbol, $\epsilon(b)$ is unchanged by permutation of the basis elements, so we may assume that $b_1 = b_1'$.
Then $b_1' \cdot b_1' = a_1$, and we can write
\begin{align*}
	\epsilon(b) &= (a_1, a_2 \cdots a_n) \prod_{2 \le i < j} (a_i, a_j) = \\
	&= (a_1, d(Q) a_1) \prod_{2 \le i < j} (a_i, a_j).
\end{align*}
Analogously, \[\epsilon(b') = (a_1', d(Q) a_1) \prod_{2 \le i < j} (a_i', a_j'). \]
The result follows by combining $a_1 = a_1'$ and the induction hypothesis applied to the orthogonal complement of $e_1$.
\end{proof}

\iffalse
\begin{lemma}\label{lem:HM_inv_pm_one}
\uses{def:HM_inv}
\lean{QuadraticForm.hasseMinkowskiInv.eq_one_or_neg_one}
\leanok
For every quadratic form $f$, we have $\epsilon(f) = \pm 1$.
\end{lemma}
\begin{proof}
Each term in the product is equal to $\pm 1$.
\end{proof}
\fi

\begin{lemma}\label{lem:represents_zero_iff_rk_3}
\lean{QuadraticForm.represents_zero_iff_of_rank_three}
\leanok
\uses{lem:HM_inv_well_defined, thm:bilinear}
Let $k$ be either $\Q_p$ or $\R$. Let $f$ be a quadratic form of rank $3$ over $k$. Then $f$ represents $0$ if and only if $(-1, -d(f)) = \epsilon(f)$.
\end{lemma}

\begin{proof}
\leanok
$f$ is equivalent to $aX^2 + bY^2 + cZ^2$ for some $a, b, c \in k^\times$, and $f$ represents $0$ if and only if the form $-caX^2 -cbY^2 - Z^2$ represents $0$.

By the definition of Hilbert symbol, this last form represents $0$ if and only if $(-ca, -cb) = 1$.
This is equivalent to
\[ (-1, -1)(-1, a)(-1, b)(c, c)(a, b)(a, c)(b, c) = 1. \]
By Proposition \ref{prop:hilbert_symbol_properties} (vi), $(c, c) = (-1, c)$, so this can be rewritten as
\[ 1 = (-1, -1)(-1, abc)(a, b)(a, c)(b, c) = (-1, -d(f))\epsilon(f), \]
or equivalently, $(-1, -d(f)) = \epsilon(f)$.
\end{proof}

\begin{lemma}\label{lem:represents_iff_rk_2}
\lean{QuadraticForm.represents_iff_of_rank_two}
\leanok
\uses{lem:represents_zero_iff_rk_3}
Let $k$ be either $\Q_p$ or $\R$. Let $f$ be a quadratic form of rank $2$ over $k$, and let $a \in k^\times / {k^\times}^2$. Then $f$ represents $a$ if and only if $(a, -d(f)) = \epsilon(f)$.
\end{lemma}
\begin{proof}
\leanok
Up to equivalence, we may assume $f = X^2 + bY^2$ for some $b \in k^\times$. By Corollary \ref{cor:represents_iff_sub_sub}, $f$ represents $a$ if and only if $f_a := f \dotminus aZ^2$ represents $0$, which by Lemma \ref{lem:represents_zero_iff_rk_3} is equivalent to $(-1, -d(f_a)) = \epsilon(f_a)$.

A computation shows that $d(f_a) = -a d(f)$ and $\epsilon(f_a) = (a, -d(f)) \epsilon (f)$. Combining this with Proposition \ref{prop:hilbert_symbol_properties}, we conclude the result.
\end{proof}

\section{Hasse--Minkowski}\label{sec:Hasse--Minkowski}

The goal of this section is to prove the Hasse--Minkowski theorem over the rational numbers. To simplify the formalization process, rather than including a single long proof of this theorem, we will split it into smaller lemmas that can be formalized independently.

Denote by $V$ the union of the set of prime numbers and the symbol $\infty$, and denote $\Q_\infty = \R$.

Let $f$ be a quadratic form over $\Q$. For each $v \in V$, the injection $\Q \to \Q_v$ allows one to view $f$ as a quadratic form over $\Q_v$, which we will denote $f_v$.

\begin{definition}\label{def:everywhere_locally_isotropic}
\lean{QuadraticForm.EverywhereLocallyIsotropic}
\uses{def:isotropic}
\leanok
A quadratic form $f$ with coefficients in $\Q$ is \emph{everywhere locally isotropic} if $f_v$ is isotropic for all $v \in V$.
\end{definition}

\begin{theorem}[Hasse--Minkowski]\label{thm:Hasse-Minkowski}
\lean{QuadraticForm.HasseMinkowski}
\uses{lem:HM_of_degenerate, lem:ev_loc_isotr_of_isotr,  lem:isotr_of_rk_zero, lem:isotr_of_rk_one, lem:isotr_of_rk_two, lem:isotr_of_rk_three, lem:isotr_of_rk_four, lem:isotr_of_five_le_rk}
\leanok
Let $f$ be a quadratic form over $\Q$. Then $f$ is isotropic if and only if $f$ is everywhere locally isotropic. In other words, $f$ has a global zero if and only if $f$ has everywhere a local zero.
\end{theorem}
\begin{proof}
Follows from the results below.
\end{proof}

\begin{theorem}[Hasse--Minkowski]\label{lem:HM_of_degenerate}
	\lean{QuadraticMap.nondegenerate_of_anisotropic}
	\uses{def:everywhere_locally_isotropic}
	\leanok
	Let $f$ be a degenerate quadratic form over $\Q$. Then $f$ is isotropic if and only if $f$ is everywhere locally isotropic. In other words, $f$ has a global zero if and only if $f$ has everywhere a local zero.
\end{theorem}
\begin{proof}
\leanok
Follows from the fact that every anisotropic quadratic form is nondegenerate.
\end{proof}

From now on, given \ref{lem:HM_of_degenerate}, we may assume that $f$ is a nondegenerate quadratic form.

\begin{lemma}\label{lem:ev_loc_isotr_of_isotr}
\lean{QuadraticMap.Isotropic.everywhereLocallyIsotropic}
\uses{def:everywhere_locally_isotropic}
\leanok
Let $f$ be an isotropic quadratic form over $\Q$. Then $f$ is everywhere locally isotropic.
\end{lemma}
\begin{proof}
\leanok
Any nontrivial solution over $Q$ yields a solution over $Q_v$ for any $v \in V$.
\end{proof}

\begin{lemma}\label{lem:isotr_of_rk_zero}
	\lean{QuadraticForm.EverywhereLocallyIsotropic.isotropic_of_rank_zero}
	\uses{def:everywhere_locally_isotropic}
	\leanok
	Let $f$ be a quadratic form over $\Q$ of rank zero which is everywhere locally isotropic. Then $f$ is isotropic.
\end{lemma}
\begin{proof}
\leanok
Suffices to see that every rank zero quadratic form over a field is anisotropic.
\end{proof}

\begin{lemma}\label{lem:equiv_sum_squares}
\lean{QuadraticForm.isotropic_iff_weightedSumSquares_units_of_nondegenerate}
\uses{def:quadratic_form}
\leanok
Let $K$ be a field of characteristic not equal to $2$, let $M$ be a $K$-vector space of positive dimension $n$ and let $f : M \to K$ be a
nondegenerate quadratic form. Then there exists a form $X_1^2 +  a_2X_2^2 + \cdots + a_nX_n^2$ for some $a_i \in K^\times$ such that it is isotropic if and only if $f$ is.
\end{lemma}
\begin{proof}
\leanok
Note that \mathlib contains a similar statement\footnote{\code{QuadraticForm.equivalent_weightedSumSquares_units_of_nondegenerate'}}, where the coefficient of $X_1$ is not assumed to be $1$.
\end{proof}

\begin{lemma}\label{lem:equiv_sum_squares_sqfree}
	\lean{QuadraticForm.isotropic_iff_weightedSumSquares_squarefree_units_of_nondegenerate}
	\uses{lem:equiv_sum_squares}
	\leanok
	Let $K$ be a field of characteristic not equal to $2$, let $M$ be a $K$-vector space of positive dimension $n$ and let $f : M \to K$ be a
	nondegenerate quadratic form. Then there exists a form $X_1^2 +  a_2X_2^2 + \cdots + a_nX_n^2$, for some squarefree $a_i \in K^\times$ such that it is isotropic if and only if $f$ is.
\end{lemma}
\begin{proof}
\leanok
Follows from Lemma \ref{lem:equiv_sum_squares}.
\end{proof}

\begin{lemma}\label{lem:isotr_of_rk_one}
	\lean{QuadraticForm.EverywhereLocallyIsotropic.isotropic_of_rank_one}
	\uses{def:everywhere_locally_isotropic}
	\leanok
	Let $f$ be a quadratic form over $\Q$ of rank one which is everywhere locally isotropic. Then $f$ is isotropic.
\end{lemma}
\begin{proof}
\leanok
Suffices to see that every nonzero quadratic form of rank one over a field is anisotropic.
\end{proof}

\begin{lemma}\label{lem:isotr_of_rk_two}
	\lean{QuadraticForm.EverywhereLocallyIsotropic.isotropic_of_rank_two}
	\uses{def:everywhere_locally_isotropic, lem:equiv_sum_squares, lem:represents_iff}
	\leanok
	Let $f$ be a quadratic form over $\Q$ of rank two which is everywhere locally isotropic. Then $f$ is isotropic.
\end{lemma}
\begin{proof}
\leanok
$f$ is equivalent to $X_1^2 - aX_2^2$, for some nonzero $a$, which is positive because $f_\infty$ represents $0$. Write
\[ a = \prod_p p^{v_p(a)}.\]
Since $f_p$ represents $0$, $a$ is a square in $\Q_p$, so $v_p(a)$ is even. It follows that $a$ is a square in $\Q$ and $f$ represents $0$.
\end{proof}

\begin{lemma}\label{lem:isotr_of_rk_three}
	\lean{QuadraticForm.EverywhereLocallyIsotropic.isotropic_of_rank_three}
	\uses{def:everywhere_locally_isotropic, lem:equiv_sum_squares_sqfree, prop:common_root, lem:represents_iff}
	\leanok
	Let $f$ be a quadratic form over $\Q$ of rank three which is everywhere locally isotropic. Then $f$ is isotropic.
\end{lemma}
\begin{proof}
$f$ is equivalent to $X_1^2 - a X_2^2 - b X_3^2$; since we can multiply $a$ and $b$ by squares, we can reduce to the case where $a$ and $b$ are square-free integers. We can also assume $|a| \le |b|$. We use induction on $m := |a| + |b|$.

If $m = 2$, then $f \sim X_1^2 \pm X_2^2 \pm X_3^2$. The case $X_1^2 + X_2^2 + X_3^2$ is excluded by the hypothesis that $f_\infty$ represents zero. In the other cases, $f$ represents zero (take e.g. $(1, 1, 0)$ or $(1, 0, 1)$, respectively).

Suppose now that $m > 2$, that is, $|b| \ge 2$ and write $b = \pm p_1 \cdots p_k$, where the $p_i$ are distinct primes. Let $p$ be one of the $p_i$; we will show that $a$ is a square modulo $p$. This is obvious if $a \equiv 0 \mod p$. Otherwise $a$ is a $p$-adic unit, and by hypothesis there exists $(x, y, z) \in \Q_p^3$ such that $z^2 - ax^2 - by^2 = 0$. By Proposition \ref{prop:common_root}, we can suppose that $(x, y, z)$ is primitive.
From $z^2 - a x^2 \equiv 0 \mod p$, we deduce that if $x \equiv 0 \mod p$, then $z \equiv 0 \mod p$ and $p^2$ divides $by^2$; since $v_p(b) = 1$, this implies $y \equiv 0 \mod p$, contradicting the assumption that $(x, y, z)$ is primitive. Therefore $x \not\equiv 0 \mod p$ and $a$ is a square modulo $p$.

Since $\Z / b\Z = \prod_i \Z / p_i \Z$, we also have that $a$ is a square modulo $b$, so there exist integers $t, b'$ such that $bb' = t^2 - a$ and we can choose $t$ so that $|t| \le |b|/2$. Then $bb'$ is a norm of the extension $k(\sqrt{a})/k$ for $k = \Q$, $\R$ or $\Q_p$. From this we conclude that $f$ represents $0$ in $k$ if and only if the same is true for $f' := X_1^2 - a X_2^2 - b' X_3^2$ (see the proof of \cite[\S III.1.1, Proposition 1]{Ser73}).

In particular, $f'$ represents $0$ in each $\Q_v$. But we have
\[ |b'| = \left | \frac{t^2 - a}{b} \right | \le \frac{|b|}{4} + 1 < |b| \quad \text{because } |b| \ge 2.  \]

Write $b' = b'' u^2$, with $b'', u$ integers and $b''$ square-free, so $|b''| < |b|$. We can then apply the induction hypothesis to the form $X_1^2 - a X_2^2 - b'' X_3^2$, which is equivalent to $f'$; hence $f'$ represents $0$ in $\Q$ and so does $f$.
\end{proof}

\begin{lemma}\label{lem:isotr_of_rk_four}
\lean{QuadraticForm.EverywhereLocallyIsotropic.isotropic_of_rank_four}
\uses{def:everywhere_locally_isotropic, lem:equiv_sum_squares, cor:prod_isotropic_iff, thm:product_formula, lem:represents_iff, thm:existence,
cor:represents_iff_sub_sub, lem:isotr_of_rk_three, lem:represents_iff_rk_2}
\leanok
Let $f$ be a quadratic form over $\Q$ of rank four which is everywhere locally isotropic. Then $f$ is isotropic.
\end{lemma}
\begin{proof}
\leanok
Up to equivalence, we may assume $f = aX_1^2 + bX_2^2 - (cX_3^2 + dX_4^2)$. Fix a place $v \in V$. Since $f_v$ represents $0$, Corollary \ref{cor:prod_isotropic_iff} shows that there exists $x_v \in \Q_v^\times$ which is represented both by $aX_1^2 + bX_2^2$ and by  $cX_3^2 + dX_4^2$.

By Lemma \ref{lem:represents_iff_rk_2}, this is equivalent to saying that
\[ (x_v, -ab)_v = (a, b)_v \quad \text{and} \quad (x_v, -cd)_v = (c, d)_v \: \text{ for all } v \in V. \]

Since $\displaystyle \prod_{v \in V} (a, b)_v = \prod_{v \in V} (c, d)_v = 1$ by Theorem \ref{thm:product_formula}, we can deduce from Theorem \ref{thm:existence} the existence of $x \in \Q^\times$ such that
\[ (x, -ab)_v = (a, b)_v \quad \text{and} \quad (x, -cd)_v = (c, d)_v \: \text{ for all } v \in V. \]

This shows that the form $aX_1^2 + bX_2^2 - xZ^2$ represents zero in each $\Q_v$, so by Lemma \ref{lem:isotr_of_rk_three}, it also represents $0$ in $\Q$. By Corollary \ref{cor:represents_iff_sub_sub}, $x$ is represented by $aX_1^2 + bX_2^2$. The same argument shows that $x$ is represented by $cX_1^2 + dX_2^2$. By applying Corollary \ref{cor:prod_isotropic_iff} again, we conclude that $f$ represents $0$.
\end{proof}

\begin{lemma}\label{lem:isotr_of_five_le_rk}
\lean{QuadraticForm.EverywhereLocallyIsotropic.isotropic_of_five_le_rank}
\uses{def:everywhere_locally_isotropic, lem:equiv_sum_squares, thm:padic_squares_open, thm:approximation, def:HM_inv, lem:isotr_of_rk_four, lem:represents_iff, def:discr}
\leanok
Let $f$ be a quadratic form over $\Q$ of rank $n \ge 5$ which is everywhere locally isotropic. Then $f$ is isotropic.
\end{lemma}
\begin{proof}
The proof uses induction on $n$. Write $f = h - g$, where $h = a_1X_1^2 + a_2X_2^2$, $g = -(a_3X_3^2 + \cdots + a_nX_n^2)$.
Let $S$ be the finite subset of $V$ consisting of $\infty, 2$, and the primes $p$ such that $v_p(a_i) \ne 0$ for $i \ge 3$.
Let $v \in S$. Since $f_v$ represents $0$, there exists $c_v \in \Q_v^\times$ which is represented in $\Q_v$ by $h$ and $g$, that is, there exist $x_i^v \in \Q_v$, $i=1, \dots, n$ (with $x_i^v \ne 0$ for $i=1, 2$, and some $i \ge 3$) such that
\[ h(x_1^v, x_2^v) = c_v = g(x_3^v, \cdots, x_n^v).\]
By  Theorem \ref{thm:padic_squares_open}, the squares of $\Q_p^\times$ form an open set. Combining this with the approximation theorem \ref{thm:approximation}, we deduce the existence of $x_1, x_2 \in \Q^\times$ such that, denoting $c := h(x_1, x_2)$, one has $c/c_v \in {\Q_v^\times}^2$ for all $v \in S$.

Consider the form $f_1 := cZ^2 - g$. For each $v \in S$, $g$ represents $c_v$ in $\Q_v$, and hence also $c$ because $c/c_v \in {\Q_v^\times}^2$, hence $f_1$ represents $0$ in $\Q_v$ by Corollary \ref{cor:represents_iff_sub_sub}. For $v \notin S$, the coefficients $-a_3, \dots, -a_n$ of $g$ are $v$-adic units, so the discriminant $d_v(g)$ is a $v$-adic unit and $\epsilon_v(g)=1$.

In all cases, $f_q$ represents $0$ in all $\Q_v$, and since the rank of $f$ is $n-1$, the inductive hypothesis shows that $f_1$ represents $0$ in $\Q$, or equivalently by Corollary \ref{cor:represents_iff_sub_sub}, $g$ represents $c$ in $\Q$. But $h$ also represents $c$, so  by Corollary \ref{cor:prod_isotropic_iff} $f$ represents $0$.

\end{proof}

